\documentclass[journal]{IEEEtran}
\usepackage[utf8]{inputenc}
\usepackage{adjustbox}
\usepackage[T1]{fontenc}
\usepackage{float}
\usepackage[spanish]{babel}
\usepackage{graphicx}

\begin{document}

\begin{titlepage}
    \centering
    \vspace*{1cm}
    \includegraphics[width=0.4\textwidth]{upiiz_logo.png}\par
    \vspace{1cm}
    {\LARGE Instituto Politecnico Nacional\par}
    \vspace{1cm}
    {\Large Unidad Profesional Interdisciplinaria de Ingeniería Campus Zacatecas IPN\par}
    \vspace{1cm}
    {\large Alumno: Francisco Javier Calderon Corrales\par}
    \vspace{0.5cm}
    {\large Trabajo: Practica 01: Analisis de casos\par}
    \vspace{0.5cm}
    {\large Carrera: Ingeniería en sistemas computacionales\par}
    \vspace{0.5cm}
    {\large Grupo: 3CM1\par}
    \vspace{0.5cm}
    {\large Docente: Erika Sánchez-Femat\par}
    \vspace{0.5cm}
    {\large Boleta: 2022670255\par}
    \vspace{1cm}
    {\large Fecha: \today\par}
\end{titlepage}

\begin{abstract}
    En este documento se abordará el tema del análisis de casos 
    para los algoritmos, más específicamente del algoritmo de 
    ordenamiento por el método de burbuja, se hablara acerca del 
    algoritmo de burbuja y su implementación, así como la versión 
    optimizada de este algoritmo, por último se detalla la complejidad 
    del algoritmo por la notación big-oh
\end{abstract}

\section{Introducción}
En el vasto universo de la informática y la programación, la optimización 
de algoritmos es una tarea esencial para mejorar el rendimiento y la 
eficiencia de las aplicaciones y procesos. Uno de los algoritmos de 
ordenamiento más básicos y fundamentales es el algoritmo de ordenamiento 
por burbuja. Sin embargo, la simplicidad de este algoritmo se traduce en 
un rendimiento deficiente cuando se trata de grandes conjuntos de datos. 
En respuesta a esta limitación, se ha desarrollado una variante conocida 
como el algoritmo de burbuja optimizada, diseñado para mitigar la ineficiencia 
de su contraparte original.

El objetivo principal de este informe es explorar y comparar estos dos 
algoritmos de clasificación: el algoritmo de burbuja clásico y su versión 
optimizada. Examinaremos sus fundamentos teóricos, su funcionamiento en 
términos de tiempo y espacio, y cómo cada uno aborda los desafíos asociados 
con la clasificación de datos. Además, realizaremos un análisis exhaustivo 
de casos de uso reales para comprender mejor en qué situaciones uno puede 
destacar sobre el otro.


\section{Desarrollo}
\subsection{implementación de los algoritmos}
Implementar los algoritmos de burbuja y burbuja optimizada fue una tarea 
sencilla, para comenzar la implementación lo primero es importar 
las librerías que ocuparemos para desarrollar el programa cumpliendo 
con los requerimientos estipulados en la práctica, utilice las librerías 
de time para determinar el tiempo de ejecución, random para que el usuario 
tenga la opción de llenar el arreglo de forma aleatoria y os para poder 
controlar la ejecución y poder pausar la ejecución
\begin{figure}[H]
    \centering
    \begin{adjustbox}{clip=false}
      \includegraphics[width=0.2\textwidth]{1.PNG}
    \end{adjustbox}
    \caption{importacion de las librerias}
  \end{figure}


Después se continúa con la declaración de las variables a 
utilizar así como las definiciones de las funciones, 
utilice una función por cada algoritmo es decir que tenemos
2 funciones una con el algoritmo de burbuja normal y otra 
con el algoritmo de burbuja optimizada, también como se 
puede apreciar agregue unos métodos que guardan el tiempo 
actual en una variable esto para saber el tiempo de ejecución.


\begin{figure}[H]
    \centering
    \begin{adjustbox}{clip=false}
      \includegraphics[width=0.5\textwidth]{2.PNG}
    \end{adjustbox}
    \caption{definicion de las funciones}
  \end{figure}

Por último se realizó el menú con las diferentes 
opciones necesarias para las pruebas y para los 
requerimientos de las prácticas.

\begin{figure}[H]
    \centering
    \begin{adjustbox}{clip=false}
      \includegraphics[width=0.5\textwidth]{3.PNG}
    \end{adjustbox}
    \caption{Una parte de la implementacion del menu}
  \end{figure}

\subsection{Analisis de casos}
Ahora ha llegado el momento de probar los algoritmos, 
para esto probaremos los dos algoritmos para el mejor 
caso, el peor caso y el caso medio, dando 3 ejemplos de 
cada uno. Primero probaremos el algoritmo de burbuja 
clásico para el mejor caso.

\begin{figure}[H]
    \centering
    \begin{adjustbox}{clip=false}
      \includegraphics[width=0.5\textwidth]{4.PNG}
    \end{adjustbox}
    \caption{Preparacion para el mejor caso de burbuja clasico}
  \end{figure}

  Como se puede apreciar en la figura 4 seleccionamos 
  el ordenamiento de burbuja y también introducimos un 
  arreglo de 1000 elementos como mejor caso, es decir 
  que el array ya está ordenado ahora es momento de ejecutarlo.

  \begin{figure}[H]
    \centering
    \begin{adjustbox}{clip=false}
      \includegraphics[width=0.5\textwidth]{5.PNG}
    \end{adjustbox}
    \caption{Ejecución del ordenamiento de burbuja clasico}
  \end{figure}

  Como se observa en la figura 5 el arreglo quedó igual debido 
  a que es el mejor caso, el algoritmo solo tuvo que comprobar 
  que todo estuviera ordenado sin realizar ningún cambio por 
  ende el tiempo de ejecución fue muy reducido.

  Ahora usaremos el algoritmo de ordenamiento por 
  burbuja optimizada, manteniendo el mejor caso y 
  el mismo tamaño de arreglo.
  \begin{figure}[H]
    \centering
    \begin{adjustbox}{clip=false}
      \includegraphics[width=0.5\textwidth]{6.PNG}
    \end{adjustbox}
    \caption{Preparacion para el mejor caso de burbuja optimizada}
  \end{figure}

  Como podemos ver en la figura 6 tenemos seleccionado
   el algoritmo de burbuja optimizada y como ya se 
   mencionó mantenemos el mejor caso y el mismo tamaño
   del arreglo, ahora ejecute el algoritmo de burbuja optimizada
   \begin{figure}[H]
    \centering
    \begin{adjustbox}{clip=false}
      \includegraphics[width=0.5\textwidth]{7.PNG}
    \end{adjustbox}
    \caption{Ejecución del ordenamiento de burbuja optimizada}
  \end{figure}
  En la figura 7 se aprecia que el tiempo de ejecución 
  fue tan bajo que tiende a cero, esto porque el 
  algoritmo de burbuja optimizada utiliza una bandera 
  para ver si en una iteración se realizó un cambio, 
  esto es para que en caso de que el algoritmo no haya 
  hecho algún cambio entonces sale del ciclo dando por 
  terminado el ordenamiento, esto aunque no haya terminado 
  el total de las iteraciones que debería.


  Ahora vamos a hacer el análisis del caso medio, 
  para esto utilizare un arreglo llenado aleatoriamente, 
  pero, manteniendo los 1000 elementos, como en el análisis 
  del mejor caso utilizare primeramente el algoritmo de 
  burbuja y posteriormente el de burbuja optimizada, es 
  importante recalcar que al ser el caso medio tendre que 
  ejecutar el algoritmo por lo menos 3 veces con diferentes 
  arreglos pero con el mismo tamaño, esto para promediar el 
  tiempo de ejecución de un caso medio para los dos algoritmos 
  de ordenamiento, ahora es momento de ejecutarlo.

  \begin{figure}[H]
    \centering
    \begin{adjustbox}{clip=false}
      \includegraphics[width=0.5\textwidth]{8.PNG}
    \end{adjustbox}
    \caption{Preparacion Para la primera ejecución de burbuja clasica}
  \end{figure}
  \begin{figure}[H]
    \centering
    \begin{adjustbox}{clip=false}
      \includegraphics[width=0.5\textwidth]{9.PNG}
    \end{adjustbox}
    \caption{Primera ejecución del ordenamiento de burbuja clasica}
  \end{figure}
  \begin{figure}[H]
    \centering
    \begin{adjustbox}{clip=false}
      \includegraphics[width=0.5\textwidth]{10.PNG}
    \end{adjustbox}
    \caption{Preparacion Para la segunda ejecución de burbuja clasica}
  \end{figure}
  \begin{figure}[H]
    \centering
    \begin{adjustbox}{clip=false}
      \includegraphics[width=0.5\textwidth]{11.PNG}
    \end{adjustbox}
    \caption{Segunda ejecución del ordenamiento de burbuja clasica}
  \end{figure}
  \begin{figure}[H]
    \centering
    \begin{adjustbox}{clip=false}
      \includegraphics[width=0.5\textwidth]{12.PNG}
    \end{adjustbox}
    \caption{Preparacion Para la tercera ejecución de burbuja clasica}
  \end{figure}
  \begin{figure}[H]
    \centering
    \begin{adjustbox}{clip=false}
      \includegraphics[width=0.5\textwidth]{13.PNG}
    \end{adjustbox}
    \caption{Tercera ejecución del ordenamiento de burbuja clasica}
  \end{figure}
  Ahora continuaré con el algoritmo de burbuja optimizada.
  \begin{figure}[H]
    \centering
    \begin{adjustbox}{clip=false}
      \includegraphics[width=0.5\textwidth]{14.PNG}
    \end{adjustbox}
    \caption{Preparacion para la primera ejecución de burbuja optimizada}
  \end{figure}
  \begin{figure}[H]
    \centering
    \begin{adjustbox}{clip=false}
      \includegraphics[width=0.5\textwidth]{15.PNG}
    \end{adjustbox}
    \caption{Primera ejecución del ordenamiento de burbuja optimizada}
  \end{figure}
  \begin{figure}[H]
    \centering
    \begin{adjustbox}{clip=false}
      \includegraphics[width=0.5\textwidth]{16.PNG}
    \end{adjustbox}
    \caption{Preparacion para la segunda ejecución de burbuja optimizada}
  \end{figure}
  \begin{figure}[H]
    \centering
    \begin{adjustbox}{clip=false}
      \includegraphics[width=0.5\textwidth]{17.PNG}
    \end{adjustbox}
    \caption{Segunda ejecución del ordenamiento de burbuja optimizada}
  \end{figure}
  \begin{figure}[H]
    \centering
    \begin{adjustbox}{clip=false}
      \includegraphics[width=0.5\textwidth]{18.PNG}
    \end{adjustbox}
    \caption{Preparacion para la tercera ejecución de burbuja optimizada}
  \end{figure}
  \begin{figure}[H]
    \centering
    \begin{adjustbox}{clip=false}
      \includegraphics[width=0.5\textwidth]{19.PNG}
    \end{adjustbox}
    \caption{Tercera ejecución del ordenamiento de burbuja optimizada}
  \end{figure}
  Como podemos apreciar no hay mucha diferencia entre los 
  tiempos de ejecución de burbuja y burbuja optimizada, 
  sorprendentemente el algoritmo de burbuja clásica tuvo 
  tiempos de ejecución ligeramente menores, es necesario 
  recalcar que los tiempos de ejecución de los algoritmos 
  depende del volumen de los datos a ordenar y también del 
  arreglo puede haber casos medios que estén más cerca de 
  ser el peor caso y otros que estén más cerca de ser el 
  mejor caso. Ahora continuaré ejecutando el peor caso 
  para los dos algoritmos.
  \begin{figure}[H]
    \centering
    \begin{adjustbox}{clip=false}
      \includegraphics[width=0.5\textwidth]{20.PNG}
    \end{adjustbox}
    \caption{Preparacion para ejecución de peor caso de burbuja clasica}
  \end{figure}
  \begin{figure}[H]
    \centering
    \begin{adjustbox}{clip=false}
      \includegraphics[width=0.5\textwidth]{21.PNG}
    \end{adjustbox}
    \caption{Ejecución de peor caso de burbuja clasica}
  \end{figure}
  \begin{figure}[H]
    \centering
    \begin{adjustbox}{clip=false}
      \includegraphics[width=0.5\textwidth]{22.PNG}
    \end{adjustbox}
    \caption{Preparacion para ejecución de peor caso de burbuja optimizada}
  \end{figure}
  \begin{figure}[H]
    \centering
    \begin{adjustbox}{clip=false}
      \includegraphics[width=0.5\textwidth]{23.PNG}
    \end{adjustbox}
    \caption{Ejecución de peor caso de burbuja optimizada}
  \end{figure}
  Aquí se puede apreciar que prácticamente no hay 
  diferencia entre el tiempo de ejecución de burbuja 
  y burbuja optimizada, esto porque en el peor caso 
  los dos algoritmos hacen prácticamente lo mismo.

  Por último clasificaré los algoritmos según la notación 
  Big-O. La notación
   Big-O (O) se utiliza para describir la complejidad 
   temporal de los algoritmos en términos de su peor 
   caso, es decir, cómo aumenta el tiempo de ejecución 
   del algoritmo en función del tamaño de la entrada (n).
   \subsection{Burbuja clasica}
   Notación Big-O (Peor Caso): $O(n^{2})$
 En el peor caso, el algoritmo de burbuja compara y 
 realiza intercambios en cada par de elementos del 
 arreglo, lo que lleva a una complejidad cuadrática 
 $(n^2)$. Esto significa que a medida que el tamaño del 
 arreglo $(n)$ aumenta, el tiempo de ejecución crece 
 cuadráticamente. Esto lo convierte en una opción 
 ineficiente para grandes conjuntos de datos, ya que el 
 tiempo de ejecución aumenta de manera significativa y 
 puede volverse prohibitivo.
 \subsection{Burbuja optimizada}
 Notación Big-O (Peor Caso): $O(n^{2})$
Justificación: A pesar de su nombre, el algoritmo 
de burbuja optimizada conserva la misma notación 
Big-O que el algoritmo de burbuja en el peor caso. 
La optimización en el algoritmo de burbuja optimizada 
radica en reducir el número de comparaciones y asignaciones 
realizadas durante la ejecución, pero aún así, en el peor 
caso, sigue teniendo una complejidad cuadrática. La razón 
de esto es que el algoritmo todavía tiene que recorrer 
todo el arreglo en cada iteración, incluso si no se realizan 
intercambios en una pasada completa.
\subsection{Comparacion de resultados}
tanto el algoritmo de burbuja como el algoritmo 
de burbuja optimizada tienen una complejidad 
temporal en el peor caso de $O(n^{2})$. La diferencia 
principal radica en la cantidad de operaciones 
realizadas en cada iteración, siendo la burbuja 
optimizada ligeramente más eficiente debido a menos 
comparaciones y asignaciones. Sin embargo, ambos 
algoritmos no son ideales para conjuntos de datos 
grandes y se recomienda su uso en situaciones donde 
el tamaño del conjunto de datos es relativamente 
pequeño o cuando se requiere un algoritmo de clasificación 
simple y fácil de implementar.

\section{Conclusion}
En este análisis comparativo de los algoritmos de 
burbuja y burbuja optimizada, hemos explorado sus 
fundamentos teóricos, su complejidad temporal y su 
rendimiento en diversos casos de uso. Estos dos 
algoritmos, a pesar de ser relativamente simples, 
ofrecen perspectivas valiosas sobre cómo se pueden 
optimizar y evaluar los algoritmos de clasificación.

En términos de su complejidad temporal en el peor caso, 
tanto el algoritmo de burbuja como su versión optimizada 
exhiben una complejidad cuadrática $O(n^{2})$, lo que significa 
que su eficiencia disminuye significativamente a medida 
que aumenta el tamaño del conjunto de datos. Sin embargo, 
en el mejor caso, ambos algoritmos operan de manera lineal 
$O(n)$ cuando el arreglo ya está ordenado, lo que destaca su 
simplicidad y eficiencia en situaciones específicas.

La principal diferencia entre estos dos algoritmos radica 
en su rendimiento en casos intermedios y su capacidad para 
manejar conjuntos de datos en los que solo se necesitan
 unos pocos intercambios para ordenarlos. La burbuja optimizada
  reduce el número de comparaciones y asignaciones en 
  comparación con la burbuja clásica, lo que la convierte 
  en una opción ligeramente más eficiente en la mayoría de 
  los casos. Sin embargo, es importante destacar que ambos 
  algoritmos tienden a mostrar un rendimiento similar en la 
  práctica, ya que el caso medio se acerca más al peor caso
   debido a la necesidad de múltiples comparaciones e 
   intercambios.

% Otras secciones y contenido

\end{document}
